\begin{table}[ht]
\centering\small
\begin{tabular}{rrrrrr}
\toprule
$T$ & $L$ & Normalized error & Error $L/T^2$ & Comparator & Ratio \\
\midrule
0.0002 & 64 & 1.1426e-07 & 182.819 & 38.062 & 4.803 \\
0.0002 & 128 & 5.7130e-08 & 182.815 & 38.062 & 4.803 \\
0.0002 & 256 & 2.8565e-08 & 182.813 & 38.062 & 4.803 \\
0.0002 & 512 & 1.4282e-08 & 182.812 & 38.062 & 4.803 \\
0.0005 & 64 & 5.8357e-07 & 149.394 & 38.351 & 3.895 \\
0.0005 & 128 & 2.9177e-07 & 149.388 & 38.351 & 3.895 \\
0.0005 & 256 & 1.4588e-07 & 149.385 & 38.351 & 3.895 \\
0.0005 & 512 & 7.2941e-08 & 149.383 & 38.351 & 3.895 \\
\bottomrule\end{tabular}
\caption{Classical normalized splitting check: $128\times128$ grid,
$p=1$, $a=0.2$, $b=0.5$, $m=4$, using the sine variant
$u_0=1+b\sin(2\pi m x)\sin(2\pi y)$. This has the same
analytic norm ratio as the cosine example. The comparator reconstructs
the printed prefactor with the initial continuum derivative norm
and the semidiscrete final norm. Finite numerical evidence supplements,
rather than proves, the analytic coefficient objection.}
\end{table}
